% !TeX program = xelatex
\documentclass[11pt]{article}
% Shared typography and author block; build with XeLaTeX, BibTeX, XeLaTeX twice.
\usepackage[a4paper,margin=25mm]{geometry}
\usepackage{amsmath,amssymb,amsthm,mathtools,bm}
\usepackage{booktabs,array,tabularx,longtable}
\usepackage{cite}
\usepackage[hidelinks,unicode]{hyperref}
\usepackage{xurl,authblk,enumitem}
\setlist[enumerate]{leftmargin=2em,itemsep=0.3ex,topsep=0.5ex}
\setlength{\emergencystretch}{3em}
\linespread{1.08}
\clubpenalty=10000
\widowpenalty=10000
\displaywidowpenalty=10000
\numberwithin{equation}{section}
\newtheorem{theorem}{Theorem}[section]
\newtheorem{lemma}[theorem]{Lemma}
\newtheorem{proposition}[theorem]{Proposition}
\newtheorem{corollary}[theorem]{Corollary}
\theoremstyle{definition}
\newtheorem{definition}[theorem]{Definition}
\theoremstyle{remark}
\newtheorem{remark}[theorem]{Remark}
\newcommand{\R}{\mathbb R}
\newcommand{\C}{\mathbb C}
\newcommand{\Repart}{\operatorname{Re}}
\newcommand{\conv}{\operatorname{conv}}
\newcommand{\SPR}{\mathrm{SPR}}
\newcommand{\code}[1]{\texttt{\small\detokenize{#1}}}
\newcommand{\id}[1]{\nolinkurl{#1}}
\newcommand{\doi}[1]{doi: \href{https://doi.org/#1}{\nolinkurl{#1}}}
\renewcommand{\Authfont}{\small}
\renewcommand{\Affilfont}{\small}
\renewcommand{\Authsep}{, }
\renewcommand{\Authand}{, }
\renewcommand{\Authands}{, }
\author[1]{Zhengtai Xiong}
\author[1]{Tiancheng Xu}
\author[1]{Xiang Ji}
\author[1]{Guowei Dou}
\author[2,3]{Long Wang}
\author[1]{Wensheng Yu}
\affil[1]{School of Electronic Engineering, Beijing University of Posts and Telecommunications, Beijing 100876, China}
\affil[2]{Center for Systems and Control, College of Engineering, Peking University, Beijing 100871, China}
\affil[3]{Center for Multi-Agent Research, Institute for Artificial Intelligence, Peking University, Beijing 100871, China}
\date{}

\title{A Lean Formalization of the SPR Synthesis Theorem for Cubic Polynomial Segments}
\hypersetup{pdftitle={A Lean Formalization of the SPR Synthesis Theorem for Cubic Polynomial Segments}}
\begin{document}
\maketitle
\begin{abstract}
Yu and Huang proved that a segment of monic cubic Hurwitz polynomials admits a common strictly positive real numerator precisely when the whole segment is Hurwitz stable. We present an independent Lean formalization of this established theorem, following its ellipse-intersection and degree-raising strategy. Explicit support functions replace the common-tangent calculation, and a completed-square identity supplies a frequency-polynomial lower bound. These two ingredients permit a direct proof of sufficiency without classifying intersections involving the triangular parts of the feasible regions. A common quadratic numerator is raised to a monic cubic numerator by an explicit scalar choice. The converse is proved through a coefficient-level numerator-stability lemma and reciprocation. The formal theorem uses the cubic Routh--Hurwitz coefficient predicate and imports only Mathlib. This work is a streamlined formalization of the original geometric method, rather than a new synthesis criterion.
\end{abstract}
\noindent\textbf{Keywords:} strictly positive real synthesis; cubic polynomial segment; ellipse intersection; Routh--Hurwitz criterion; Lean.

\section{Introduction and relation to the original proof}
For a family of uncertain denominators, robust strictly positive real (SPR) synthesis asks for a single numerator that makes every resulting transfer function SPR. Even for two stable endpoint polynomials, stability of the endpoints alone is insufficient: the entire segment must remain stable. Yu and Huang established the necessary and sufficient segment condition for degrees at most three \cite{yuHuang1999}. Their cubic proof identifies an ellipse and a triangular region in the coefficient plane of quadratic numerators, reduces common feasibility to intersection of two ellipses, and raises the numerator degree by perturbation.

The present paper describes the complete theorem proved in \code{SPR/simplified/n3.lean}. Its central geometric object is exactly the ellipse in Lemma~1 of \cite{yuHuang1999}; the intersection argument implements the separating-line idea of their Lemma~6. The construction $a+Rq$ is a positive rescaling of $q+ra$ from their Lemma~5, with $r=1/R$. Thus this construction is not a different underlying method. The changes are explicit support-function algebra and a uniform lower bound for the real-part numerator, which together remove the need for the triangular-region cases in the sufficiency proof.

The organization follows the mathematical presentation of the supplied arbitrary-order manuscript \cite{referenceDraft}: the synthesis statement is followed by its frequency reduction, the central exclusion argument, and the final construction. The cubic proof itself does not use the contact-certificate method of that manuscript. Nor does the Lean file import an arbitrary-order synthesis theorem. A comparison with the original proof and an exact statement of the verification boundary are given below.

\section{Definitions and main theorem}
Write
\[
 a(s)=s^3+a_1s^2+a_2s+a_3,\qquad
 b(s)=s^3+b_1s^2+b_2s+b_3,
 \qquad p_\lambda=(1-\lambda)a+\lambda b.
\]
In this paper the coefficient predicate $H_3(a)$ means
\begin{equation}\label{eq:rh}
 a_1>0,\quad a_2>0,\quad a_3>0,\quad a_1a_2>a_3.
\end{equation}
These are the classical cubic Routh--Hurwitz inequalities. They characterize roots in the open left half-plane. The Lean source takes \eqref{eq:rh} as its definition; it does not separately formalize the general Routh--Hurwitz equivalence with root location. All formal stability assertions below refer to this explicit predicate.

For a real cubic numerator $c$ with positive leading coefficient, $c/a$ is called SPR when $H_3(a)$ holds and
\begin{equation}\label{eq:spr}
 \Repart\frac{c(i\omega)}{a(i\omega)}>0\qquad(\omega\in\R).
\end{equation}
Condition \eqref{eq:rh} implies $a(i\omega)\ne0$: the real and imaginary equations for a nonzero imaginary-axis root would give $\omega^2=a_2$ and $a_3=a_1a_2$. At zero frequency, $a_3>0$ excludes a root. Monic equal-degree quotients tend to $1$ at infinity. On the remaining compact frequency interval, continuity then gives a positive minimum of the real part. Quadratic numerators used during the construction are described only as having strict frequency positivity; they are not called equal-degree SPR numerators.

\begin{theorem}[Yu--Huang criterion, coefficient form]\label{thm:main}
Suppose $H_3(a)$ and $H_3(b)$. The following conditions are equivalent:
\begin{enumerate}
\item $H_3(p_\lambda)$ for every $\lambda\in[0,1]$;
\item there is a monic cubic $c$ such that $c/a$ and $c/b$ are SPR;
\item there is a cubic $c$ with positive leading coefficient such that $c/a$ and $c/b$ are SPR.
\end{enumerate}
Under these conditions there also exists a monic quadratic $q$ with strictly positive real parts of $q/a$ and $q/b$ at every finite frequency.
\end{theorem}

The equivalence of the two cubic normalizations follows by division by the positive leading coefficient. The rest of the paper proves the first two conditions equivalent and constructs the quadratic witness.

\section{Elliptic regions and a positivity margin}
For $q(s)=s^2+xs+y$ put $t=\omega^2$ and define
\begin{align}
 F_a(t)&=(a_1-x)t^2+(a_2x-a_1y-a_3)t+a_3y,\label{eq:F}\\
 E_a(x,y)&=(a_2x-a_1y-a_3)^2-4(a_1-x)a_3y,\label{eq:E}\\
 \Omega_a&=\{(x,y)\in\R^2:E_a(x,y)<0\}.
\end{align}
Direct multiplication yields
\begin{equation}\label{eq:freq}
 \Repart\frac{q(i\omega)}{a(i\omega)}
 =\frac{F_a(\omega^2)}{|a(i\omega)|^2}.
\end{equation}
Expansion of \eqref{eq:E} gives the conic in \cite[Lemma~1]{yuHuang1999}:
\[
 a_2^2x^2+2(2a_3-a_1a_2)xy+a_1^2y^2
 -2a_2a_3x-2a_1a_3y+a_3^2.
\]

\begin{lemma}\label{lem:margin}
If $H_3(a)$ and $(x,y)\in\Omega_a$, then $a_1-x>0$ and
\begin{equation}\label{eq:margin}
 F_a(t)\ge m_a:=\frac{-E_a(x,y)}{4(a_1-x)}>0\qquad(t\in\R).
\end{equation}
\end{lemma}
\begin{proof}
Set $\Delta_a=a_1a_2-a_3>0$. Translating to the center gives
\[
 E_a(a_1/2+X,a_2/2+Y)
 =a_2^2X^2+2(2a_3-a_1a_2)XY+a_1^2Y^2-a_3\Delta_a.
\]
Multiplication of the quadratic part by $a_1^2$ completes a square:
\[
 a_1^2\bigl(E_a(x,y)+a_3\Delta_a\bigr)
 =\bigl(a_1^2Y+(2a_3-a_1a_2)X\bigr)^2
 +4a_3\Delta_a X^2.
\]
Consequently $E_a(x,y)<0$ implies $|X|<a_1/2$, hence $0<x<a_1$. Finally,
\[
 4(a_1-x)F_a(t)
 =\bigl(2(a_1-x)t+a_2x-a_1y-a_3\bigr)^2-E_a(x,y).
\]
Division by $4(a_1-x)>0$ proves \eqref{eq:margin}.
\end{proof}

This is a lower bound for the polynomial $F_a$, not a uniform positive lower bound for $\Repart(q/a)$: the latter tends to zero as frequency tends to infinity. Retaining this distinction is essential in the degree-raising argument.

\section{Stable segments force ellipse intersection}
For a nonzero normal $n=(A,B)$ define
\begin{equation}\label{eq:dual}
 D_a(A,B)=a_1^2A^2-2(2a_3-a_1a_2)AB+a_2^2B^2,
 \qquad r_a=\sqrt{D_a(A,B)}.
\end{equation}
The matrix of the translated conic is positive definite, with determinant $4a_3\Delta_a$. Its inverse gives the exact support values
\begin{equation}\label{eq:support}
 \sup_{(x,y)\in\Omega_a}(Ax+By)
 =\frac{a_1A+a_2B+r_a}{2},\qquad
 \inf_{(x,y)\in\Omega_a}(Ax+By)
 =\frac{a_1A+a_2B-r_a}{2}.
\end{equation}
The extrema here are suprema and infima of an open ellipse. They are approached by interior points along the segment from the center to a boundary support point. This is the limiting argument implemented in Lean; attainment in the open set is not assumed. Positive definiteness also gives $r_a>0$.

\begin{lemma}\label{lem:intersection}
If $H_3(p_\lambda)$ for all $\lambda\in[0,1]$, then $\Omega_a\cap\Omega_b\ne\varnothing$.
\end{lemma}
\begin{proof}
Suppose the two nonempty open convex ellipses are disjoint. Convex separation supplies a nonzero normal, oriented so that \eqref{eq:support} implies
\[
 r_a+r_b\le L:=(b_1-a_1)A+(b_2-a_2)B.
\]
The dual quadratic has the affine-segment identity
\begin{equation}\label{eq:dualblend}
 D_{p_\lambda}(A,B)
 =(1-\lambda)D_a(A,B)+\lambda D_b(A,B)
 -\lambda(1-\lambda)L^2.
\end{equation}
Take $\lambda_*=r_a/(r_a+r_b)\in(0,1)$. Since
\[
 (1-\lambda_*)r_a^2+\lambda_*r_b^2=r_ar_b
 =\lambda_*(1-\lambda_*)(r_a+r_b)^2,
\]
\eqref{eq:dualblend} gives $D_{p_{\lambda_*}}(A,B)\le0$. This contradicts positive definiteness of the dual form for the Hurwitz polynomial $p_{\lambda_*}$.
\end{proof}

The contradiction is the separating-line mechanism of \cite[Lemma~6]{yuHuang1999}, expressed without choosing axis intercepts or assuming their differences are nonzero. It also handles identical endpoint polynomials: the separation hypothesis then immediately fails.

\section{Construction of an equal-degree numerator}
Choose $(x,y)\in\Omega_a\cap\Omega_b$ and let $q=s^2+xs+y$. Lemma~\ref{lem:margin} proves strict frequency positivity of both quadratic quotients. For the cubic construction, introduce the endpoint cross polynomial
\begin{align}
 G_{a,b}(t)={}&t^3+(a_1b_1-a_2-b_2)t^2\notag\\
 &+(a_2b_2-a_3b_1-a_1b_3)t+a_3b_3.
\end{align}
It satisfies $G_{a,b}(\omega^2)=\Repart(a(i\omega)\overline{b(i\omega)})$. Being monic of odd degree, it is bounded below on $[0,\infty)$. More explicitly, if $G=t^3+At^2+Bt+C$, set
\[
 M=|A|+|B|+|C|+1,\qquad
 K=4|A|M^2+2|B|M+|C|+1.
\]
For $0\le t<2M$, termwise bounds show $G(t)>-K$; for $t\ge2M$, the cubic term dominates the three lower terms and gives $G(t)\ge0$. Hence the following choice is valid:
\begin{equation}\label{eq:R}
 R=\frac{K+1}{m_b}>0,\qquad
 c=a+Rq=s^3+(a_1+R)s^2+(a_2+Rx)s+(a_3+Ry).
\end{equation}
Indeed,
\[
 \Repart(c/a)=1+R\Repart(q/a)>0,
\]
whereas
\[
 \Repart\frac{c(i\omega)}{b(i\omega)}
 =\frac{G_{a,b}(\omega^2)+R F_b(\omega^2)}{|b(i\omega)|^2}>0
\]
because the numerator is greater than $-K+Rm_b=1$. This proves sufficiency in Theorem~\ref{thm:main}.

\section{Necessity and the whole segment}
The formal proof does not assume that an SPR numerator is stable. It proves the needed cubic statement algebraically.
\begin{lemma}\label{lem:numstable}
If $H_3(a)$ and a monic cubic $c=s^3+c_1s^2+c_2s+c_3$ satisfies \eqref{eq:spr}, then $H_3(c)$.
\end{lemma}
\begin{proof}
Let $r=a_3/a_1$, so $0<r<a_2$. At $t=0$, positivity of the cross polynomial gives $c_3>0$. At $t=r$, the real part of the denominator vanishes and its odd part has positive sign, so cross positivity gives $c_2>r>0$. At $t=a_2$, the odd part of the denominator vanishes and its real part is negative; positivity therefore implies $c_1a_2>c_3$, and $c_1>0$.

Set $u=c_3/c_1>0$. If $c_2\le u$, evaluation at $t=u$ eliminates the real part of $c$ and gives
\[
 u(c_2-u)(a_2-u)>0.
\]
Equality $c_2=u$ is impossible, and the strict inequality would force $u>a_2$, contradicting $c_1a_2>c_3$. Thus $c_2>u$ and $c_1c_2>c_3$.
\end{proof}

Suppose $c/a$ and $c/b$ are SPR. Their numerators cannot vanish on the imaginary axis, and
\[
 \Repart(z/w)>0\quad\Longleftrightarrow\quad\Repart(w/z)>0
 \qquad(z,w\ne0).
\]
Lemma~\ref{lem:numstable} makes $c$ a stable denominator. The identity
\[
 \Repart\frac{p_\lambda(i\omega)}{c(i\omega)}
 =(1-\lambda)\Repart\frac{a(i\omega)}{c(i\omega)}
 +\lambda\Repart\frac{b(i\omega)}{c(i\omega)}>0
\]
and another application of Lemma~\ref{lem:numstable} give $H_3(p_\lambda)$. This proves necessity.

\begin{corollary}
The numerator in Theorem~\ref{thm:main} works for every denominator on the segment.
\end{corollary}
\begin{proof}
The real cross product with $c$ is affine in the denominator coefficients and is positive at both endpoints. It remains positive at each convex combination. The denominator is stable by necessity, so division by its squared modulus proves the claim.
\end{proof}

\section{Formal verification}
The proof presented in this paper has been formally verified in Lean. The code repository is \url{https://github.com/XPiPiT/n_random/tree/main/SPR/simplified}.

\section{Conclusion}
The cubic synthesis criterion is formally established with the same ellipse-intersection and degree-raising strategy as Yu and Huang. Explicit dual support values make the separation contradiction uniform, while a polynomial margin yields an explicit degree-raising scale. The resulting formalization proves both directions of the established theorem and isolates the reduced-degree intermediate result without reproducing the original triangular-region classification.

\bibliographystyle{IEEEtran}
\bibliography{spr_references}
\end{document}
